\centerline{\bf \Huge 10 класс}

\problem{10.1}
Разбейте какой-нибудь клетчатый квадрат на клетчатые квадратики так, чтобы не все квадратики были одинаковы, но квадратиков каждого размера было одно и то же количество.

\problem{10.2}
Петя и Вася по очереди выписывают на доску натуральные числа, не превосходящие 2018 (выписывать уже имеющееся число запрещено); начинает Петя. Если после хода игрока на доске оказываются три числа, образующих арифметическую прогрессию, - этот игрок выигрывает. У кого из игроков есть стратегия, позволяющая ему гарантированно выиграть?

\problem{10.3}
Положительные числа $x, y$ таковы, что $x^5-y^3 \geqslant 2 x$. Докажите, что $x^3 \geqslant 2 y$.

\problem{10.4}
Пусть $O$ - центр окружности $\Omega$, описанной около остроугольного треугольника $A B C$. На дуге $A C$ этой окружности, не содержащей точку $B$, взята точка $P$. На отрезке $B C$ выбрана точка $X$ так, что $P X \perp A C$. Докажите, что центр окружности, описанной около треугольника $B X P$, лежит на окружности, описанной около треугольника $A B O$.
    
\problem{10.5}
Дано нечётное число $n>10$. Найдите количество способов расставить по кругу в некотором порядке натуральные числа $1,2,3, \ldots, n$ так, чтобы каждое число являлось делителем суммы двух соседних с ним чисел. (Способы, отличающиеся поворотом или отражением, считаются одинаковыми.)